\subsection{Product Measures and Fubini Theory}

% **Role:** modern integration and almost-everywhere analysis.
%
% Product measures, Tonelli, Fubini, iterated integration, applications.

Product measure builds measurable rectangles and reads higher-dimensional integrals through slices.

\begin{definition} \label{def:lebesgue-product-slices} \hlt{Products and Slices}\\
For $A\subseteq\R^n$, $B\subseteq\R^k$, set $A\times B\subseteq\R^{n+k}$.\\
For $E\subseteq\R^n\times\R^k$ define $E_x\defeq\{y\in\R^k:(x,y)\in E\}$ as \hlt{$x$-slice}, and  $E^y\defeq\{x\in\R^n:(x,y)\in E\}$ as \hlt{$y$-slice}.\\
For $f:E\to\R$, write $f_x(y)=f^y(x)\defeq f(x,y)$ where defined.
\end{definition}

\begin{theorem} \label{thm:measurable-products-product-measure} \hlt{Measurable Products}\\
If $A\subseteq\R^n$, $B\subseteq\R^k$ Lebesgue measurable, then $A\times B$ measurable, $m(A\times B)=m(A)m(B)$; $0\cdot\infty=\infty\cdot0=0$.
\begin{proof}
For boxes $I,J$, $m(I\times J)=m(I)m(J)$; disjoint box decompositions extend this to open sets.\\
First assume $A,B$ are bounded with finite measure.\\
By regularity choose closed $F_A\subseteq A$, $F_B\subseteq B$ and open $G_A\supseteq A$, $G_B\supseteq B$ such that $m(G_A\setminus F_A),\ m(G_B\setminus F_B)<\varepsilon$.\\
Then $F_A\times F_B\subseteq A\times B\subseteq G_A\times G_B$, and the open/closed product formula squeezes $m(A\times B)$ to $m(A)m(B)$.\\
For general $A,B$, set $A_N=A\cap[-N,N]^n$ and $B_N=B\cap[-N,N]^k$. Since $A_N\times B_N\uparrow A\times B$, continuity from below gives $m(A\times B)=\lim_N m(A_N\times B_N)=\lim_N m(A_N)m(B_N)=m(A)m(B)$ with $0\cdot\infty=\infty\cdot0=0$.\\
\hlt{Structural Commentary}: Product measure comes from boxes and regularity.\\
\hlt{Key Insight}: Regularity turns rough measurable products into box products plus zero-set errors.
\end{proof}
\end{theorem}

\begin{theorem} \label{thm:zero-slices-null-set-criterion} \hlt{Zero Slices}\\
If $E\subseteq\R^n\times\R^k$ is Lebesgue measurable, then $E$ is null iff $E_x$ is null for a.e. $x$.
\begin{proof}
\begin{enumerate}[label=\roman*.]
\setlength{\itemsep}{0pt}
\item[$(\Rightarrow)$] If $m(E)=0$, cover $E$ by an open $W$ with small measure.\\
For $X_\alpha(W)=\{x:m(W_x)>\alpha\}$, $m(W)\geq \alpha m(X_\alpha(W))$; letting $m(W)\to0$ gives $m(E_x)=0$ for a.e. $x$.
\item[$(\Leftarrow)$] Remove the exceptional null base so every slice is null. For compact $K\subseteq E$, cover each $K_x$ by a small-measure slice neighborhood; compactness gives finitely many product neighborhoods covering $K$, whose total measure is arbitrarily small by Theorem \ref{thm:measurable-products-product-measure}. Thus $m(K)=0$, hence $m(E)=0$ by regularity.
\end{enumerate}
\hlt{Structural Commentary}: Nullity can be checked fiberwise after removing an exceptional base set.\\
\hlt{Key Insight}: Product measure turns small slice covers into small total measure.
\end{proof}
\end{theorem}

\begin{figure}[H]
  \centering
  \begin{tikzpicture}[
  scale=0.86,
  >=Stealth,
  every node/.style={font=\scriptsize},
  axis/.style={->, black!40, line width=0.5pt},
  boundary/.style={black!75, line width=0.8pt},
  guide/.style={black!22, line width=0.45pt},
  slice/.style={blue!70!black, line width=1pt},
  labelbox/.style={fill=white, inner sep=1pt}
]
  \def\setE{(-2.85,0.45)
    .. controls (-2.55,1.22) and (-1.28,1.44) .. (-0.08,1.40)
    .. controls (1.35,1.36) and (2.78,1.05) .. (3.05,0.20)
    .. controls (3.35,-0.75) and (2.20,-1.45) .. (0.58,-1.55)
    .. controls (-1.10,-1.66) and (-2.55,-1.18) .. (-2.95,-0.28)
    .. controls (-3.08,0.02) and (-3.02,0.25) .. (-2.85,0.45)
    -- cycle}

  \draw[axis] (-3.35,-1.85) -- (3.55,-1.85) node[right] {$x$};
  \draw[axis] (-3.15,-2.05) -- (-3.15,1.85) node[above] {$y$};

  \fill[blue!12] \setE;
  \begin{scope}
    \clip \setE;
    \foreach \x in {-2.2,-1.2,-0.2,0.8,1.8} {
      \draw[guide] (\x,-1.65) -- (\x,1.55);
    }
    \draw[slice] (2.20,-1.23) -- (2.20,0.86);
  \end{scope}
  \draw[boundary] \setE;

  \draw[black!35, line width=0.45pt] (2.20,-1.85) -- (2.20,-1.98);
  \node[labelbox] at (2.20,-2.17) {$x_0$};

  \draw[blue!70!black, line width=0.55pt] (2.20,0.86) -- (2.60,1.12);
  \node[labelbox, text=blue!70!black, anchor=west] at (2.62,1.12) {$E_{x_0}$};
  \node[labelbox] at (-0.30,-1.05) {$E$};
\end{tikzpicture}

  \caption{A planar set can be read by vertical slices; Fubini integrates the slice data over the base variable.}
  \label{fig:cavalieri-slicing-planar-set}
\end{figure}

\begin{theorem} \label{thm:cavalieri-principle-euclidean-slices} \hlt{Cavalieri's Principle for Lebesgue Measure}\\
If $E\in\mathcal{L}^{n+k}$, then $E_x\in\mathcal{L}^k$ for a.e. $x$, and the a.e.-defined function $f_E(x)\defeq m_k(E_x)$ has a measurable representative with $m_{n+k}(E)=\int_{\R^n} f_E\,dm_n$. The analogous statement holds for $E^y$.
\begin{proof}
Let $P(E)$ denote the displayed assertion. For $E=A\times B$, $f_E=m_k(B)\mathbf{1}_A$, so $P(E)$ is Theorem \ref{thm:measurable-products-product-measure}.\\
Fix a bounded box $Q=I\times J$. The class $\{E\subseteq Q:P(E)\}$ contains finite unions of rectangles and is monotone: if $E_j\uparrow E$, then $f_{E_j}\uparrow f_E$ and monotone convergence applies; if $E_j\downarrow E$, then $f_{E_j}\downarrow f_E\leq m_k(J)\mathbf{1}_I$ and dominated convergence applies. Hence the monotone class theorem gives $P(E)$ for every Borel $E\subseteq Q$. Exhaustion by $Q_N=[-N,N]^{n+k}$ and monotone convergence give all Borel $E$.\\
For Lebesgue $E$, choose Borel $B$ with $m_{n+k}(E\triangle B)=0$ (Remark \ref{rem:borel-null-regular-measurable-sets}). Then $m_k((E\triangle B)_x)=0$ for a.e. $x$ (Theorem \ref{thm:zero-slices-null-set-criterion}), so $f_E=f_B$ a.e. and $m_{n+k}(E)=m_{n+k}(B)$. The $y$-slice statement is identical.\\
\hlt{Structural Commentary}: Slicing turns a measure in $\R^{n+k}$ into a measurable family of measures in $\R^k$.\\
\hlt{Key Insight}: Prove boxes, close under monotone limits, then ignore null symmetric differences.
\end{proof}
\end{theorem}

\begin{perspective} \label{per:fubini-proof-ladder} \hlt{Fubini Proof Ladder}\\
Prove Fubini first for indicators, extend to simple functions, pass to nonnegative functions by monotone convergence, then split $f=f^+-f^-$ under absolute integrability.\\
\hlt{Key Insight}: Fubini is a slice theorem, proof runs through indicators, simple functions, monotone convergence.
\end{perspective}

\begin{proposition} \label{prop:undergraph-measurability-level-sets} \hlt{Undergraphs and Level-Set Measurability}\\
Let $f:\R^n\to[0,\infty]$, set $F_a\defeq\{x:f(x)>a\}$, and write $\mathcal{U}f=\{(x,t)\in\R^n\times[0,\infty):0<t<f(x)\}$.\\
Then $f$ is Lebesgue measurable iff $F_q\in\mathcal{L}^n \ \forall q\in\Q_{>0}$. In that case, $\mathcal{U}f\in\mathcal{L}^{n+1}$ and $(\mathcal{U}f)^t=F_t$, $t>0$.
\begin{proof}
For $a\geq0$, $F_a=\bigcup\limits_{\substack{q\in\Q\\q>a}}F_q$; for $a<0$, $F_a=\R^n$.\\
Hence rational positive thresholds are equivalent to measurability.\\
If they are measurable, then $\mathcal{U}f =\bigcup\limits_{\substack{q,r\in\Q\\0<q<r}} F_r\times(q,r)\in\mathcal{L}^{n+1}$. Finally, $(\mathcal{U}f)^t=\{x:0<t<f(x)\}=F_t$ for $t>0$.\\
\hlt{Structural Commentary}: A nonnegative function may be studied as the region lying beneath its graph.\\
\hlt{Key Insight}: Horizontal slices of an undergraph are exactly strict superlevel sets.
\end{proof}
\end{proposition}

\begin{theorem} \label{thm:fubini-tonelli-nonnegative-euclidean} \hlt{Fubini--Tonelli for Nonnegative Functions}\\
Let $f:\R^n\times\R^k\to[0,\infty]$ be Lebesgue measurable. Then $f_x$ is measurable for a.e. $x$; with $F(x)\defeq\int_{\R^k}f_x\,dm_k$ there and arbitrary values elsewhere, $F$ is measurable and $\int_{\R^{n+k}} f\,dm_{n+k}=\int_{\R^n}F\,dm_n$.\\
The same holds for $G(y)\defeq\int_{\R^n}f^y\,dm_n$, with $\int_{\R^{n+k}}f\,dm_{n+k}=\int_{\R^k}G\,dm_k$.
\begin{proof}
Put $U\defeq\mathcal{U}f\subseteq\R^{n+k}\times[0,\infty)$. Proposition \ref{prop:undergraph-measurability-level-sets} gives $U\in\mathcal{L}^{n+k+1}$. By Cavalieri, Theorem \ref{thm:cavalieri-principle-euclidean-slices}, $U_x$ is measurable for a.e. $x$, $H(x)\defeq m_{k+1}(U_x)$ is measurable, and $m_{n+k+1}(U)=\int_{\R^n}H\,dm_n$.\\
For those $x$, $U_x=\mathcal{U}(f_x)$, so $f_x$ is measurable and $H(x)=\int_{\R^k}f_x\,dm_k$.\\
Also $m_{n+k+1}(U)=\int_{\R^{n+k}}f\,dm_{n+k}$ by the undergraph formula. The $y$-formula is identical.\\
\hlt{Structural Commentary}: Tonelli's theorem is Cavalieri's principle applied one dimension higher.\\
\hlt{Key Insight}: $(\mathcal{U}f)_x=\mathcal{U}(f_x)$.
\end{proof}
\end{theorem}

\begin{theorem} \label{thm:fubini-absolutely-integrable-functions} \hlt{Fubini for $L^1$ Functions}\\
If $f\in L^1(\R^2)$, then $f(x,\cdot)\in L^1(\R)$ for a.e. $x$ and $f(\cdot,y)\in L^1(\R)$ for a.e. $y$. Defining $F(x)\defeq\int_{\R}f(x,y)\,dy$; $G(y)\defeq\int_{\R}f(x,y)\,dx$ on those sets and as $0$ elsewhere gives $F,G\in L^1(\R)$ and $\int_{\R}F(x)\,dx=\int_{\R^2}f=\int_{\R}G(y)\,dy$
\begin{proof}
Let $E_x=\{y:(x,y)\in E\}$ and $E^y=\{x:(x,y)\in E\}$. For measurable $E\subseteq\R^2$, the box/regularity argument and Theorems \ref{thm:measurable-products-product-measure}--\ref{thm:zero-slices-null-set-criterion} give $\int_{\R}m(E_x)\,dx=m(E)=\int_{\R}m(E^y)\,dy$.\\
Hence Fubini holds for $\mathbf{1}_E$, and by linearity for every nonnegative simple function.\\
Let $0\leq h$ be measurable and choose $0\leq s_n\uparrow h$ by Theorem \ref{thm:simple-approximation-measurable-functions}.\\
Monotone convergence gives $\int_{\R}\!\left(\int_{\R}h(x,y)\,dy\right)dx =\lim_n\int_{\R}\!\left(\int_{\R}s_n(x,y)\,dy\right)dx =\lim_n\int_{\R^2}s_n =\int_{\R^2}h$ and the same argument with $x,y$ interchanged. Now let $f\in L^1(\R^2)$.\\
Applying the nonnegative case to $\abs{f}$ yields $\int_{\R}\!\left(\int_{\R}\abs{f(x,y)}\,dy\right)dx =\int_{\R^2}\abs{f}<\infty$, so $f(x,\cdot)\in L^1(\R)$ for a.e. $x$; similarly $f(\cdot,y)\in L^1(\R)$ for a.e. $y$. With $f=f^+-f^-$, define, up to null-set choices, $F=F^+-F^-$, $F^\pm(x)=\int_{\R}f^\pm(x,y)\,dy$, and analogously $G=G^+-G^-$. Then $\int_{\R}\abs{F}\leq\int_{\R}(F^++F^-)=\int_{\R^2}\abs{f}<\infty$, and similarly $G\in L^1(\R)$. Subtracting the nonnegative identities for $f^+$ and $f^-$ gives $\int_{\R}F=\int_{\R^2}f=\int_{\R}G$.\\
\hlt{Structural Commentary}: Prove identity for $\mathbf{1}_E$, extend by linearity and monotone convergence, apply to $f^\pm$.\\
\hlt{Key Insight}: $\int\abs{f}<\infty$ makes the positive and negative slice integrals finite a.e., so subtraction is legitimate.
\end{proof}
\end{theorem}
